\documentclass[12pt, a4paper]{ctexart}

% ── 页面设置 ──────────────────────────────────────────────
\usepackage[top=3cm, bottom=2.5cm, left=3cm, right=2.5cm, headheight=14pt]{geometry}

% ── 字体与排版 ────────────────────────────────────────────
\usepackage{fontspec}
\usepackage{amsmath, amssymb, amsthm}
\usepackage{bm}

% ── 图表 ──────────────────────────────────────────────────
\usepackage{pgfplots}
\pgfplotsset{compat=1.18}
\usepackage{tikz}
\usetikzlibrary{shapes.geometric, arrows.meta, positioning, calc, decorations.pathreplacing}
\usepackage{graphicx}
\usepackage{caption}
\usepackage{subcaption}
\usepackage{float}
\captionsetup{font=small, labelfont=bf}

% ── 代码 ──────────────────────────────────────────────────
\usepackage{listings}
\usepackage{xcolor}

\definecolor{codebg}{HTML}{F8F8F8}
\definecolor{codekey}{HTML}{0033B3}
\definecolor{codecmt}{HTML}{8C8C8C}
\definecolor{codestr}{HTML}{067D17}
\definecolor{codenum}{HTML}{1750EB}

\lstset{
  language=C++,
  basicstyle=\small\ttfamily,
  keywordstyle=\color{codekey}\bfseries,
  commentstyle=\color{codecmt}\itshape,
  stringstyle=\color{codestr},
  numberstyle=\tiny\color{gray},
  backgroundcolor=\color{codebg},
  frame=single,
  framerule=0.4pt,
  rulecolor=\color{gray!40},
  numbers=left,
  numbersep=8pt,
  tabsize=4,
  breaklines=true,
  showstringspaces=false,
  xleftmargin=1.5em,
}

% ── 页眉页脚 ──────────────────────────────────────────────
\usepackage{fancyhdr}
\pagestyle{fancy}
\fancyhf{}
\fancyhead[L]{\small\color{gray}数据结构与算法设计 2 \; 期中大作业}
\fancyhead[R]{\small\color{gray}2451951 \; 于沛可}
\fancyfoot[C]{\thepage}
\renewcommand{\headrulewidth}{0.4pt}

% ── 超链接与目录 ──────────────────────────────────────────
\usepackage[colorlinks=true, linkcolor=black, urlcolor=blue!70!black]{hyperref}
\usepackage{booktabs}
\usepackage{enumitem}

% ── 自定义颜色 ────────────────────────────────────────────
\definecolor{heapcolor}{HTML}{D62728}
\definecolor{randcolor}{HTML}{1F77B4}
\definecolor{lincolor}{HTML}{2CA02C}
\definecolor{tjiblue}{HTML}{003087}
\definecolor{cellT}{HTML}{AED6F1}
\definecolor{cellF}{HTML}{F0F0F0}
\definecolor{mstedge}{HTML}{E74C3C}

% ── 定理环境 ──────────────────────────────────────────────
\newtheorem{theorem}{定理}
\newtheorem{proposition}{命题}

% ──────────────────────────────────────────────────────────
\begin{document}

% ═══════════════════════════════════════════════════════════
%  封面
% ═══════════════════════════════════════════════════════════
\thispagestyle{empty}
\begin{center}
  \vspace*{2cm}
  {\zihao{3}\heiti 同济大学}\\[0.4em]
  \vspace{1.5cm}

  \noindent\rule{\linewidth}{1.2pt}\\[0.6em]
  {\zihao{2}\heiti 数据结构与算法设计 2\\[0.4em] 期中大作业报告}\\[0.6em]
  \noindent\rule{\linewidth}{1.2pt}

  \vspace{1.5cm}

  {\zihao{-3}\heiti 分治、动态规划与贪心算法}\\[0.3em]
  {\zihao{4}\songti 第k小元素 \; $\cdot$ \; 卡牌均分 \; $\cdot$ \; 最小生成树}

  \vspace{2.5cm}

  \begin{tabular}{rl}
    \zihao{4}\heiti 学\quad 号 & \zihao{4}\songti 2451951 \\[0.8em]
    \zihao{4}\heiti 姓\quad 名 & \zihao{4}\songti 于沛可 \\[0.8em]
    \zihao{4}\heiti 专\quad 业 & \zihao{4}\songti 数据科学与大数据技术 \\[0.8em]
    \zihao{4}\heiti 课\quad 程 & \zihao{4}\songti 数据结构与算法设计 2 \\[0.8em]
    \zihao{4}\heiti 提交日期 & \zihao{4}\songti 2026 年 5 月
  \end{tabular}
\end{center}
\newpage

% ═══════════════════════════════════════════════════════════
%  目录
% ═══════════════════════════════════════════════════════════
\tableofcontents
\newpage

% ═══════════════════════════════════════════════════════════
%  编程环境
% ═══════════════════════════════════════════════════════════
\section{编程语言与实验环境}

编程语言采用 C++（C++11 标准），编译器为 g++ 10.3.0（TDM-GCC 64-bit），
运行平台为 Windows 10（x86-64）。
时间测量使用 \texttt{std::chrono::high\_resolution\_clock}，单位为微秒（$\mu$s）。
功能验证采用 \texttt{-O2} 编译；性能基准测试采用 \texttt{-O0}，
以避免编译器消除死代码对结果的干扰。

\newpage

% ═══════════════════════════════════════════════════════════
%  第一题
% ═══════════════════════════════════════════════════════════
\section{第一题 \; 分治算法：第 $k$ 小元素}

\subsection{问题描述}

给定无序整数序列 $A[1..n]$ 及整数 $k$（$1 \le k \le n$），
求第 $k$ 小元素，即排序后位于第 $k$ 位的值。

\subsection{三种算法}

\subsubsection{方法一：基于堆的选择}

维护容量为 $k$ 的大根堆 $H$：将前 $k$ 个元素入堆后，
对每个后续元素 $a$，若 $a < \mathrm{top}(H)$，则弹出堆顶并压入 $a$；
遍历结束时 $\mathrm{top}(H)$ 即为答案。

\paragraph{时间复杂度} 建堆 $O(k)$；后续 $n-k$ 次操作各需 $O(\log k)$，
总计
\[
  T_1(n, k) = O\!\left(k + (n-k)\log k\right).
\]
当 $k \ll n$ 时，$\log k$ 较小，整体接近线性；
当 $k = n/2$ 时，$\log k = \log(n/2)$，复杂度退化为 $O(n\log n)$。

\subsubsection{方法二：随机划分线性选择（QuickSelect）}

借鉴快速排序的分区思想，每次随机选取 pivot 并以此将 $A[l..r]$
划分为三部分，根据 pivot 排名与 $k$ 的关系仅对相关一侧递归。

\paragraph{时间复杂度}
期望时间 $O(n)$。
若用随机 pivot，则每步期望至少减少 $1/4$ 的元素，
递推可得 $\mathbf{E}[T_2(n)] = O(n)$。
最坏情形（每次选到极端元素）为 $O(n^2)$，但对均匀随机 pivot，
其概率随 $n$ 指数衰减，工程实践中几乎不出现。

\subsubsection{方法三：中位数的中位数线性选择}

将 $A[l..r]$ 每 5 个分组，各组内排序取中位数，
再递归求这些中位数的中位数作为 pivot，
从而保证每次划分后两侧规模各不超过 $\lceil 7n/10 \rceil + 6$。

\paragraph{时间复杂度} 设 $T_3(n)$ 为最坏情形代价，则：
\begin{align*}
  T_3(n) &\le T_3(n/5) + T_3(7n/10 + 6) + O(n).
\end{align*}
归纳可证 $T_3(n) = O(n)$，且对任意输入均成立，
即这是一个\textbf{最坏情形} $O(n)$ 的算法。

\subsection{代码实现（核心逻辑）}

\begin{lstlisting}[caption={方法一：基于大根堆的选择}]
int heapSelect(vector<int> arr, int k) {
    priority_queue<int> pq;
    for (int i = 0; i < k; i++) pq.push(arr[i]);
    for (int i = k; i < (int)arr.size(); i++) {
        if (arr[i] < pq.top()) {
            pq.pop();
            pq.push(arr[i]);
        }
    }
    return pq.top();
}
\end{lstlisting}

\begin{lstlisting}[caption={方法二：随机划分选择（QuickSelect）}]
int randSelect(vector<int>& a, int l, int r, int k) {
    if (l == r) return a[l];
    int p = l + rand() % (r - l + 1);
    swap(a[p], a[r]);
    int pivot = a[r], i = l - 1;
    for (int j = l; j < r; j++)
        if (a[j] <= pivot) swap(a[++i], a[j]);
    swap(a[++i], a[r]);
    int cnt = i - l + 1;
    if (k == cnt) return a[i];
    return k < cnt ? randSelect(a, l, i-1, k)
                   : randSelect(a, i+1, r, k-cnt);
}
\end{lstlisting}

\begin{lstlisting}[caption={方法三：中位数的中位数线性选择（递归核心）}]
int medianSelect(vector<int>& a, int l, int r, int k) {
    if (r - l + 1 <= 5) {
        sort(a.begin()+l, a.begin()+r+1);
        return a[l+k-1];
    }
    vector<int> meds;
    for (int i = l; i <= r; i += 5) {
        int j = min(i+4, r);
        sort(a.begin()+i, a.begin()+j+1);
        meds.push_back(a[(i+j)/2]);
    }
    int mm = meds.size();
    int pivot = medianSelect(meds, 0, mm-1, (mm+1)/2);
    // 以 pivot 为基准划分，再对相应侧递归
    int pivotPos = l;
    for (int i = l; i <= r; i++)
        if (a[i] == pivot) { pivotPos = i; break; }
    swap(a[pivotPos], a[r]);
    int p = l-1;
    for (int j = l; j < r; j++)
        if (a[j] <= pivot) swap(a[++p], a[j]);
    swap(a[++p], a[r]);
    int rank = p - l + 1;
    if (k == rank) return a[p];
    return k < rank ? medianSelect(a, l, p-1, k)
                    : medianSelect(a, p+1, r, k-rank);
}
\end{lstlisting}

\subsection{实验结果}

\paragraph{功能验证}

对测试数组 $\{3,1,4,1,5,9,2,6,5,3,5\}$，三种方法对所有 $k\in[1,11]$
均给出正确结果，与排序后直接索引一致。

\paragraph{性能测试}

以下两组实验分别考察运行时间随 $n$（固定 $k=n/2$）和随 $k$（固定 $n=200{,}000$）的变化。
每组实验生成同一随机数组，分别传副本给三种算法。

\begin{figure}[H]
  \centering
  \begin{tikzpicture}
    \begin{axis}[
      width=0.88\linewidth, height=7cm,
      xlabel={数组规模 $n$（万）},
      ylabel={运行时间（ms）},
      xmin=0, xmax=110,
      ymin=0, ymax=560,
      xtick={5,10,50,100},
      xticklabels={5,10,50,100},
      grid=both,
      grid style={dotted, gray!30},
      legend style={at={(0.05,0.95)}, anchor=north west, font=\small},
      tick label style={font=\small},
      label style={font=\small},
      title style={font=\small\bfseries},
      title={三种算法运行时间随数组规模 $n$ 的变化（$k = \lfloor n/2 \rfloor$）},
    ]
    % 堆选择
    \addplot[color=heapcolor, mark=square*, thick, mark size=3pt]
      coordinates {(5,20.17) (10,41.62) (50,231.10) (100,508.13)};
    \addlegendentry{堆选择 $O(n\log k)$}
    % 随机选择
    \addplot[color=randcolor, mark=triangle*, thick, mark size=3.5pt]
      coordinates {(5,4.63) (50,25.0) (100,47.43)};
    \addlegendentry{随机选择 $O(n)$ 期望}
    % 线性选择
    \addplot[color=lincolor, mark=*, thick, mark size=3pt]
      coordinates {(5,8.34) (10,33.65) (50,144.14) (100,282.20)};
    \addlegendentry{线性选择（中位数）$O(n)$ 最坏}
    \end{axis}
  \end{tikzpicture}
  \caption{运行时间随 $n$ 的变化（$k = \lfloor n/2 \rfloor$，\texttt{-O0} 编译，单位 ms）}
  \label{fig:kth-n}
\end{figure}

\begin{figure}[H]
  \centering
  \begin{tikzpicture}
    \begin{axis}[
      width=0.88\linewidth, height=7cm,
      xlabel={$k$ 的取值},
      ylabel={运行时间（ms）},
      xmin=0, xmax=210000,
      ymin=0, ymax=95,
      xtick={0,50000,100000,150000,200000},
      xticklabels={0,$5\!\times\!10^4$,$10^5$,$1.5\!\times\!10^5$,$2\!\times\!10^5$},
      grid=both,
      grid style={dotted, gray!30},
      legend style={at={(0.5,0.97)}, anchor=north, font=\small},
      tick label style={font=\small},
      label style={font=\small},
      title style={font=\small\bfseries},
      title={三种算法运行时间随 $k$ 的变化（$n = 200{,}000$）},
    ]
    % 堆选择
    \addplot[color=heapcolor, mark=square*, thick, mark size=3pt]
      coordinates {(2000,9.19) (20000,41.85) (100000,83.31)
                   (180000,50.11) (200000,25.05)};
    \addlegendentry{堆选择}
    % 随机选择
    \addplot[color=randcolor, mark=triangle*, thick, mark size=3.5pt]
      coordinates {(1,7.25) (2000,8.36) (20000,8.22)
                   (100000,15.95) (200000,8.33)};
    \addlegendentry{随机选择}
    % 线性选择
    \addplot[color=lincolor, mark=*, thick, mark size=3pt]
      coordinates {(1,50.02) (2000,50.19) (20000,57.69)
                   (100000,59.03) (180000,58.05) (200000,58.31)};
    \addlegendentry{线性选择（中位数）}
    \end{axis}
  \end{tikzpicture}
  \caption{运行时间随 $k$ 的变化（$n = 200{,}000$，\texttt{-O0} 编译，单位 ms）}
  \label{fig:kth-k}
\end{figure}

\subsection{分析与讨论}

\paragraph{三种方法的实际对比}

图~\ref{fig:kth-n} 中，$n=100$ 万时随机选择（47 ms）约为中位数算法（282 ms）的 1/6，
尽管二者理论复杂度同为 $O(n)$。
常数因子的差异源于中位数算法需额外执行约 $n/5$ 规模的递归
以及多轮分组排序，使每元素的平均操作数远高于 QuickSelect。
因此，\textbf{最坏情形的渐近保证与平均实际性能之间存在显著权衡}：
中位数算法适用于输入可能被对抗性构造的场景，
而 QuickSelect 在随机数据下更具工程竞争力。

\paragraph{堆选择关于 $k$ 的单峰特性}

图~\ref{fig:kth-k} 中，堆选择的曲线在 $k = n/2$ 附近达到峰值，
两端（$k \ll n$ 或 $k \approx n$）均显著更快。
这与复杂度公式 $O(k + (n-k)\log k)$ 完全一致：
当 $k$ 极小时，$\log k \approx 0$，堆操作代价趋近于零；
当 $k \to n$ 时，$(n-k) \to 0$，同样减少操作量。

\paragraph{线性选择对 $k$ 的不敏感性}

线性选择在不同 $k$ 值下的运行时间几乎恒定在 58 ms 左右，
说明其常数因子主要由分组、提取中位数等与 $k$ 无关的步骤主导。
这一特性在 $k = 1$（求最小值）时最为显眼：
此时堆选择只需 $O(n)$ 的单次扫描（$\log 1 = 0$），
而中位数算法仍执行完整的 $O(n)$ 递归流程，代价相差约 7 倍。

\newpage

% ═══════════════════════════════════════════════════════════
%  第二题
% ═══════════════════════════════════════════════════════════
\section{第二题 \; 动态规划：宝可梦卡牌均分}

\subsection{问题描述}

给定 $n$ 张卡牌，第 $i$ 张价值为正整数 $v_i$，总价值 $S = \sum_{i=1}^{n} v_i$。
将所有卡牌划分为两个不相交的子集，使两者价值之差最小。
输出两堆各自的价值及所含卡牌编号（1-indexed）。

\subsection{方法一：暴力枚举}

\paragraph{算法描述}
枚举 $\{0,1\}^n$ 中所有二进制掩码 $m$，
第 $i$ 位为 1 表示卡牌 $i$ 归入第一堆。
计算 $s = \sum_{i:\, m_i=1} v_i$，差值为 $|S - 2s|$，取差值最小的 $m$。

\paragraph{时间复杂度} $O(2^n \cdot n)$。当 $n \ge 30$ 时约需 $10^{10}$ 次操作，实际不可行。

\subsection{方法二：动态规划}

\paragraph{问题转化}
最小差值等价于在所有价值子集的和中，
找到最接近 $\lfloor S/2 \rfloor$ 的可达和 $s^*$，
则最优差值为 $S - 2s^*$。
这是 0-1 背包的可行性变体。

\paragraph{状态定义}
令 $\textit{dp}[i][j] = \text{true}$ 当且仅当从前 $i$ 张卡牌中，
存在某个子集的价值恰为 $j$。

\paragraph{递推关系}
\begin{equation}
  \textit{dp}[i][j] =
  \textit{dp}[i-1][j]
  \;\lor\;
  \bigl(j \ge v_i \;\land\; \textit{dp}[i-1][j-v_i]\bigr),
  \quad i \in [1,n],\; j \in [0, S].
  \label{eq:dp}
\end{equation}

边界条件：$\textit{dp}[0][0] = \text{true}$，其余为 $\text{false}$。

\paragraph{求解最优分配}
填表后，令 $s^* = \max\{j \le \lfloor S/2 \rfloor : \textit{dp}[n][j] = \text{true}\}$，
以逆序回溯（若 $\textit{dp}[i-1][\cdot]$ 不能达到当前和，则第 $i$ 张必在此堆中）
还原具体分配。

\paragraph{时间与空间复杂度} $O(n \cdot S)$（时间）；$O(n \cdot S)$（存储二维表以供回溯）。
若仅需最优值，可将二维数组压缩为一维（滚动更新），空间降至 $O(S)$。

\paragraph{DP 表格可视化}

图~\ref{fig:dp-table} 展示了样例（$n=8$，$\mathbf{v}=(2,1,3,1,5,2,3,4)$，$S=21$，目标 $\lfloor S/2 \rfloor = 10$）
中各行 $\textit{dp}[i][j]$（$j \in [0, 10]$）的分布。
蓝色格表示可达，灰色格表示不可达。
关键观察：第 5 行（加入价值为 5 的卡牌后）首次使 $j=10$ 变为可达，
此后目标和始终保持可达，回溯在此行截止。

\begin{figure}[H]
\centering
\begin{tikzpicture}[scale=0.82]
  % 数据：行i, 列j=0..10 可达性
  % i=0: 只有j=0
  % i=1 v=2: 0,2
  % i=2 v=1: 0,1,2,3
  % i=3 v=3: 0..6
  % i=4 v=1: 0..7
  % i=5 v=5: 0..10 (全部)

  \def\rows{6}  % i=0..5
  \def\cols{11} % j=0..10

  % 背景格
  \foreach \i in {0,...,5} {
    \foreach \j in {0,...,10} {
      \draw[gray!30, thin] (\j, -\i) rectangle (\j+1, -\i+1);
    }
  }

  % 可达格（填蓝色）
  % i=0
  \fill[cellT] (0,0) rectangle (1,1);
  % i=1 (v=2)
  \fill[cellT] (0,-1) rectangle (1,0);
  \fill[cellT] (2,-1) rectangle (3,0);
  % i=2 (v=1)
  \foreach \j in {0,1,2,3} {
    \fill[cellT] (\j,-2) rectangle (\j+1,-1);
  }
  % i=3 (v=3)
  \foreach \j in {0,...,6} {
    \fill[cellT] (\j,-3) rectangle (\j+1,-2);
  }
  % i=4 (v=1)
  \foreach \j in {0,...,7} {
    \fill[cellT] (\j,-4) rectangle (\j+1,-3);
  }
  % i=5 (v=5) 全部可达
  \foreach \j in {0,...,10} {
    \fill[cellT] (\j,-5) rectangle (\j+1,-4);
  }

  % 重画格线（盖在填色上方）
  \foreach \i in {0,...,5} {
    \foreach \j in {0,...,10} {
      \draw[gray!50, thin] (\j, -\i) rectangle (\j+1, -\i+1);
    }
  }

  % j=10 一列特殊高亮（目标列）
  \draw[mstedge, very thick] (10, -5) rectangle (11, 1);

  % 列标签（j 值）
  \foreach \j in {0,...,10} {
    \node[above, font=\scriptsize] at (\j+0.5, 1) {$\j$};
  }
  % 行标签
  \node[left, font=\scriptsize, align=right] at (0, 0.5)  {$i=0$};
  \node[left, font=\scriptsize, align=right] at (0,-0.5)  {$i=1,\,v_1\!=\!2$};
  \node[left, font=\scriptsize, align=right] at (0,-1.5)  {$i=2,\,v_2\!=\!1$};
  \node[left, font=\scriptsize, align=right] at (0,-2.5)  {$i=3,\,v_3\!=\!3$};
  \node[left, font=\scriptsize, align=right] at (0,-3.5)  {$i=4,\,v_4\!=\!1$};
  \node[left, font=\scriptsize, align=right] at (0,-4.5)  {$i=5,\,v_5\!=\!5$\;$\star$};

  % 顶部轴标
  \node[above, font=\small\bfseries] at (5.5, 1.5) {$j$（目标和）};
  % 右侧注释
  \node[right, font=\scriptsize, text=mstedge] at (11.1, -4.5)
    {$j\!=\!10$ 首次可达};
\end{tikzpicture}
\caption{DP 状态表（$i=0\sim5$，$j=0\sim10$）。
  蓝色：\textit{dp}[i][j]=true；灰色：false。
  红框标注目标列 $j=10$，行 $i=5$ 是首个使 $j=10$ 可达的行。}
\label{fig:dp-table}
\end{figure}

\subsection{代码实现（核心逻辑）}

\begin{lstlisting}[caption={动态规划求解（含回溯）}]
void dpSolve(vector<int>& v) {
    int n = v.size(), total = 0;
    for (int x : v) total += x;
    int half = total / 2;

    // 建立 (n+1) x (total+1) 的 dp 表
    vector<vector<bool>> dp(n+1, vector<bool>(total+1, false));
    dp[0][0] = true;
    for (int i = 1; i <= n; i++)
        for (int j = 0; j <= total; j++) {
            dp[i][j] = dp[i-1][j];
            if (j >= v[i-1] && dp[i-1][j-v[i-1]])
                dp[i][j] = true;
        }

    // 找最大可达和 <= half
    int best = 0;
    for (int j = half; j >= 0; j--)
        if (dp[n][j]) { best = j; break; }

    // 逆序回溯还原第一堆
    vector<int> pile1;
    int cur = best;
    for (int i = n; i >= 1; i--)
        if (!dp[i-1][cur]) {
            pile1.push_back(i);
            cur -= v[i-1];
        }
    sort(pile1.begin(), pile1.end());
    // ... 输出部分略
}
\end{lstlisting}

\subsection{实验结果}

\paragraph{功能验证}

样例 $\mathbf{v}=(2,1,3,1,5,2,3,4)$（$n=8$）：两种方法均输出
$s_1=10$（卡牌 1、3、5）、$s_2=11$（卡牌 2、4、6、7、8），差值为 1，为最优解。

\paragraph{暴力与 DP 运行时间对比}

\begin{figure}[H]
  \centering
  \begin{tikzpicture}
    \begin{axis}[
      ybar,
      width=0.7\linewidth, height=6.5cm,
      ymode=log,
      log basis y=10,
      ylabel={运行时间（ms，对数轴）},
      xlabel={数组规模 $n$},
      symbolic x coords={$n=16$, $n=20$},
      xtick=data,
      ymin=0.0001, ymax=1000,
      bar width=18pt,
      enlarge x limits=0.4,
      grid=major,
      grid style={dotted, gray!30},
      legend style={at={(0.97,0.97)}, anchor=north east, font=\small},
      tick label style={font=\small},
      label style={font=\small},
      title style={font=\small\bfseries},
      title={暴力枚举与动态规划运行时间对比（对数纵轴）},
      nodes near coords,
      nodes near coords style={font=\tiny, anchor=south},
      every node near coord/.append style={yshift=2pt},
    ]
    \addplot[fill=heapcolor!80, draw=heapcolor!80!black]
      coordinates {($n=16$, 9.28) ($n=20$, 153.65)};
    \addlegendentry{暴力枚举 $O(2^n)$}
    \addplot[fill=randcolor!80, draw=randcolor!80!black]
      coordinates {($n=16$, 0.001) ($n=20$, 0.001)};
    \addlegendentry{动态规划 $O(n \cdot S)$（$\approx0$ ms）}
    \end{axis}
  \end{tikzpicture}
  \caption{两种算法运行时间对比。纵轴取对数；DP 运行时间低于计时分辨率（$<1\,\mu$s），
    以 $0.001\,\text{ms}$ 标注。}
  \label{fig:dp-compare}
\end{figure}

\subsection{分析与讨论}

图~\ref{fig:dp-compare} 直观呈现了指数级与伪多项式时间的差距：
$n$ 从 16 增至 20 时，暴力枚举耗时从 9.3 ms 增至 153.6 ms（约 $2^4 \approx 16$ 倍），
完全符合 $O(2^n)$ 的增长规律；
而 DP 始终在 $1\,\mu$s 以内，即便 $n=500$（$v_i \in [1,20]$，$S \approx 5000$）
也仅需约 10 ms。

DP 的效率来自对子问题的复用：
式~\eqref{eq:dp} 的转移结构使得每个"能否凑出和 $j$"
的子问题只需计算一次，避免了暴力法中对相同子集的反复遍历。
DP 的代价在于空间：完整二维表占用 $O(n \cdot S)$ 字节，
当 $S$ 较大时可改用一维滚动数组以换取空间效率。

\newpage

% ═══════════════════════════════════════════════════════════
%  第三题
% ═══════════════════════════════════════════════════════════
\section{第三题 \; 贪心算法：供水管网最小生成树}

\subsection{问题建模}

将 $n$ 栋楼视为图 $G=(V, E)$ 的顶点，每条潜在管道对应一条边，
权值为铺设长度。供水网络的三条要求（全连通、无环、总长最短）
恰好对应图论中的\textbf{最小生成树（MST）}。

\subsection{算法：Kruskal}

\paragraph{贪心策略}
按边权从小到大贪心地加入边，当且仅当该边不与已选边集形成环路。

\paragraph{正确性}
MST 的\emph{割性质}保证了此策略的全局最优性：
对图的任意割（将 $V$ 分为两个非空子集），
穿越该割的最小权边必属于某棵 MST。
Kruskal 每次选取当前最短的跨分量边，
恰好对各割均满足此条件，因此所选边集构成 MST。

\paragraph{环路检测}
使用带路径压缩与按秩合并的\textbf{并查集（Union-Find）}：
若边 $(u,v)$ 两端点属于同一连通分量（$\mathrm{find}(u) = \mathrm{find}(v)$），
则加入该边会形成环，跳过；否则合并两分量。

\paragraph{时间复杂度}
边排序 $O(m \log m)$；$m$ 次并查集操作均摊 $O(m\, \alpha(n))$，
其中 $\alpha$ 为反阿克曼函数，增长极缓慢，可视为常数。
总体：$O(m \log m)$。

\subsection{代码实现}

\begin{lstlisting}[caption={并查集 + Kruskal 主逻辑}]
struct Edge { int u, v, w; };

int par[1001], rnk[1001];
void init(int n) {
    for (int i = 1; i <= n; i++) { par[i] = i; rnk[i] = 0; }
}
int find(int x) {
    return par[x] == x ? x : (par[x] = find(par[x]));
}
bool unite(int x, int y) {
    x = find(x); y = find(y);
    if (x == y) return false;
    if (rnk[x] < rnk[y]) swap(x, y);
    par[y] = x;
    if (rnk[x] == rnk[y]) rnk[x]++;
    return true;
}

// Kruskal 主算法
sort(edges.begin(), edges.end(),
     [](const Edge& a, const Edge& b){ return a.w < b.w; });
init(n);
int total = 0, cnt = 0;
for (auto& e : edges) {
    if (unite(e.u, e.v)) {
        mst.push_back(e);
        total += e.w;
        if (++cnt == n-1) break;
    }
}
\end{lstlisting}

\subsection{测试用例与结果}

\paragraph{测试图}

6 栋楼，10 条候选管道，具体数据见图~\ref{fig:mst}。
带权图中，红色粗线为最小生成树所选管道，灰色虚线为排除在外的候选管道。

\begin{figure}[H]
\centering
\begin{tikzpicture}[
  node distance=2.5cm,
  every node/.style={circle, draw, fill=white, minimum size=8mm, font=\small\bfseries},
  mst/.style={draw=mstedge, very thick, line width=1.8pt},
  nst/.style={draw=gray!50, dashed, thin},
  lbl/.style={font=\scriptsize, fill=white, inner sep=1.5pt,
              text=gray!70, draw=none},
  mstlbl/.style={font=\scriptsize\bfseries, fill=white, inner sep=1.5pt,
                 text=mstedge, draw=none},
]

\node (1) at (0, 3)   {1};
\node (2) at (3, 4.5) {2};
\node (3) at (2, 1.5) {3};
\node (4) at (5, 1.5) {4};
\node (5) at (6, 4)   {5};
\node (6) at (7.5, 1.5) {6};

% 非MST边（灰色虚线）
\draw[nst] (1) -- node[lbl,above left] {4} (2);
\draw[nst] (2) -- node[lbl,above left] {8} (4);
\draw[nst] (3) -- node[lbl,above] {10} (5);
\draw[nst] (4) -- node[lbl,right] {7} (5);
\draw[nst] (5) -- node[lbl,right] {9} (6);

% MST边（红色粗线）
\draw[mst] (1) -- node[mstlbl,left] {2} (3);
\draw[mst] (2) -- node[mstlbl,above] {3} (3);
\draw[mst] (2) -- node[mstlbl,above right] {5} (5);
\draw[mst] (3) -- node[mstlbl,below] {6} (4);
\draw[mst] (4) -- node[mstlbl,below] {1} (6);

% 图例（手动用 tikz 线条画）
\draw[mst] (8.0, 4) -- (9.0, 4);
\node[draw=none, fill=none, right, font=\scriptsize] at (9.1, 4) {MST 管道};
\draw[nst] (8.0, 3.4) -- (9.0, 3.4);
\node[draw=none, fill=none, right, font=\scriptsize, text=gray!70] at (9.1, 3.4) {候选管道};

\end{tikzpicture}
\caption{测试用例图示（6 栋楼，10 条候选管道）。
  红色粗线：Kruskal 选出的最小生成树（总长 $1+2+3+5+6=\mathbf{17}$）；
  灰色虚线：未选边。}
\label{fig:mst}
\end{figure}

\paragraph{Kruskal 执行过程}

\begin{enumerate}[leftmargin=2em, itemsep=2pt]
  \item 按权值排序：$(4\text{-}6,1)$、$(1\text{-}3,2)$、$(2\text{-}3,3)$、$(1\text{-}2,4)$、
        $(2\text{-}5,5)$、$(3\text{-}4,6)$、$(4\text{-}5,7)$、$(2\text{-}4,8)$、
        $(5\text{-}6,9)$、$(3\text{-}5,10)$。
  \item $(4\text{-}6,1)$：两端不连通，\textbf{加入}，分量：$\{4,6\}$。
  \item $(1\text{-}3,2)$：\textbf{加入}，分量：$\{1,3\}$、$\{4,6\}$。
  \item $(2\text{-}3,3)$：\textbf{加入}，分量：$\{1,2,3\}$、$\{4,6\}$、$\{5\}$。
  \item $(1\text{-}2,4)$：1 与 2 同属 $\{1,2,3\}$，\textbf{跳过}（成环）。
  \item $(2\text{-}5,5)$：\textbf{加入}，分量：$\{1,2,3,5\}$、$\{4,6\}$。
  \item $(3\text{-}4,6)$：\textbf{加入}，所有节点并入同一分量，\textbf{完成}。
\end{enumerate}

最终总长 $1+2+3+5+6 = \mathbf{17}$，与程序输出一致。

\subsection{分析与讨论}

Kruskal 的贪心正确性依赖于割性质，该性质的成立条件是边权的可比较性与图的连通性。
对于本题（同材质水管、权值为长度），条件均满足。
此外，当候选管道数量 $m$ 远小于 $n^2$ 时（现实场景中地形限制使边稀疏），
Kruskal 的 $O(m \log m)$ 比 Prim 的朴素实现 $O(n^2)$ 更具优势；
若图为稠密图，则 Prim 配合 Fibonacci 堆（$O(m + n\log n)$）更优。

\newpage

% ═══════════════════════════════════════════════════════════
%  其他说明
% ═══════════════════════════════════════════════════════════
\section{其他说明}

\paragraph{问题 1：线性时间选择实现的一个细节}
中位数算法需要在原数组中找到 pivot 值并进行划分。
若数组中存在重复值，线性扫描得到的"pivot 位置"可能不是
参与分组中位数计算的那个具体元素，
但由于我们只关心 pivot 的\emph{值}而非其身份，
选取任意一个等于 pivot 值的元素参与划分均不影响正确性。

\paragraph{问题 2：计时精度的干扰}
在 Windows 上，\texttt{chrono::high\_resolution\_clock} 的实际分辨率约为 1 ms，
导致部分小规模测试出现读数为零的情况。
性能实验因此选取足够大的 $n$ 以确保数据可信；
对于理论上应有正值但实测为零的数据点，在图表中予以略去或注明。
此外，\texttt{-O2} 优化可能将返回值未被使用的函数调用整体消除，
性能测试最终使用 \texttt{volatile} 标记接收返回值以规避此问题。

\paragraph{问题 3：DP 的空间与回溯的取舍}
DP 求解卡牌分堆时，若仅需最优差值，一维滚动数组即可；
但为还原具体分配方案，需保留完整二维表以供回溯。
当 $n$ 或各 $v_i$ 较大时，内存开销可能显著。
一种折中方案是分段存储，将问题拆分为若干子区间分别处理，
但实现复杂度较高，本次作业未予实现。

\paragraph{体会}
三道题分别展示了分治、动态规划与贪心的核心设计思想。
尤为值得关注的是算法理论与工程实践之间的落差：
中位数算法在最坏情形上具有理论优越性，
但其常数因子的劣势在随机输入下被充分放大，
这说明对算法的评估不能止步于渐近符号，
实际常数、缓存行为及分支预测均会对性能产生不可忽视的影响。

\end{document}
